\section{From scientific assistance to a testable banking research loop}
\label{sec:exec24-scientific-method}

An automated research loop needs a scientific question that survives every transformation.
A servicing analyst may ask whether an earlier contact changes ten-day cash receipts,
while a text assistant extracts an eventual repayment flag and an optimizer maximizes
that flag. Each component can operate correctly while the composition answers a
different question. This is the practical value of examining the whole scientific
cycle, rather than treating a better hypothesis generator as a complete causal system.
Zhang and colleagues organize their perspective around observation, hypothesis,
experimentation and automation, with human scientific goals and evaluation as essential
constraints. We use their May 2025 preprint as a primary source for that argument
and its caveats. Its descriptions of other systems are leads to their originals,
not independent qualification of those systems.~\citep{exec24_scimethod_original}

The source is the 45-page version of \emph{Advancing the Scientific Method with Large
Language Models: From Hypothesis to Discovery}, arXiv 2505.16477v1, dated 22 May 2025.
A later journal record has a different title and the same thirteen authors. The
journal full text was not available in this inspection; neither exact equivalence
nor journal-specific claims are assumed. This distinction matters because a review
may change its qualifications or examples between versions. The preprint's account
supports a design vocabulary and criticism, not an empirical estimate of bank
productivity or a theorem that autonomous discovery will succeed.

\subsection{Observation is part of the estimand}

The perspective treats observation as a reduction of complex information, including
annotation, classification and extraction. It also acknowledges information loss when
domain data are converted into a representation accepted by a model. In banking the
representation is consequential: a loss of timing, cohort membership or assignment
information can make an otherwise answerable intervention question unanswerable.
Let $H$ denote the original account history and let $Z=g(H)$ be the assistant's
extracted representation. For a target $q(H)$, a deterministic reconstruction from
$Z$ exists on the declared history class only if $q$ is constant on every fibre
of $g$. Indeed, if $g(h_1)=g(h_2)$ but $q(h_1)\ne q(h_2)$, any proposed function
$f$ would have to give two values to $f(g(h_1))$. Conversely, constancy allows
$f(z)$ to be defined by the common target value on its fibre. This elementary
statement concerns information preservation, not stochastic causal identification.

Consider two stipulated histories of a USD 10,000 obligation: full repayment on
day 5 and full repayment on day 25. The extraction of eventual repayment is one
in both histories, while the ten-day cash target is USD 10,000 in the first and
zero in the second. An accurate eventual-repayment classifier therefore does not
resolve the liquidity target. Keeping raw histories, event times, availability
times and extraction provenance is a substantive alternative to committing early
to a convenient label. It costs more storage and review effort, but allows later
questions to be assessed for information sufficiency. Neither retained detail nor
a high extraction score alone establishes truthful assignment or absence of
unmeasured confounding.~\citep{exec24_scimethod_reference}

A business-to-causal translation should consequently produce a structured question:
population, intervention and comparator, outcome, horizon, unit, information available
at decision time, and loss function. A generated diagram then proposes mechanisms
consistent with those definitions. Its edges need separate provenance: a substantive
restriction, a temporal impossibility, an inspected empirical result, or an
unresolved hypothesis. Textual plausibility cannot decide between observationally
equivalent mechanisms. Where the question changes after seeing results, that change
belongs in a new target and evaluation record, rather than silently changing the
meaning of the original request.

\subsection{Hypothesis generation and verification have different jobs}

The source explicitly discusses a possible creative role for speculative generation,
while requiring reliability in experimental work. This is useful only when an
unsupported suggestion remains visibly hypothetical. A proposed edge from contact
to repayment delay is a candidate for comparison with alternatives,
not a fact to insert into every downstream adjustment set. Competing stories may
include selection by latent repayment capacity, contact triggered by early warning
signals, and a genuine contact effect. Their distinguishing implications must be
derived before choosing an observation or intervention that could separate them.
Retrieval can improve access to documented premises; it cannot turn a similar
population, treatment or observation clock into the same one.

The perspective's discussion of formal verification contains an especially important
qualification: translating a generated hypothesis into a formal language can itself
introduce error. For example, the goal of cost strictly below USD 100 is
$c<100$. An assistant translating it as $c\leq100$ may produce a perfectly verified
solution at $c=100$ that violates the original goal. A proof checker is entitled to
accept the weaker formal specification. The missing obligation is correspondence
between business intent and formal premises. Our exact boundary witness does not
test a proof assistant; it explains why a verified artifact must bind the agreed
question, units, inequalities, domain and assumptions as well as the solver output.
\citep{exec24_scimethod_original,exec24_scimethod_reference}

There are therefore several distinct checks. A parser can establish syntactic
validity; a type checker can reject a unit or shape mismatch; an exact small oracle
can test a proposed algorithm on a declared family; a proof can establish a theorem
under explicit premises; and an experiment can challenge those premises in a stated
population. Failure at any layer is informative, and passing one layer does not
eliminate the others. Generated executable code can falsify a claim on a witness
without proving it for all permitted inputs. A successful formal derivation also
does not make a bank's observational assignment mechanism exogenous.

\subsection{Repeated agreement is not calibrated scientific confidence}

The preprint surveys repetition, self-correction, debate and external feedback, but
also reports limitations of reasoning, planning and self-explanation. It calls for
system-level trust assessment that includes tasks, tools, workflows and feedback.
Its proposed analogy with statistical confidence is a research direction; it does
not supply a generally valid confidence interval for a causal answer produced by
an arbitrary agent. We should retain that question rather than converting it into
a claimed product feature.~\citep{exec24_scimethod_original}

An exact calculation illustrates the issue. Suppose five judges each have error
probability $p=0.1$ on a stipulated request law. Under independent errors the
majority-error probability is
\[
 \sum_{k=3}^{5}\binom{5}{k}p^k(1-p)^{5-k}
 =\frac{107}{12500}=0.00856.
\]
If every judge instead shares the same Bernoulli error, each still has marginal
error $0.1$, but majority error remains $0.1$. The shared failure could arise from
the same incorrect source, an identical omitted confounder or a common mistranslation.
The example does not estimate any real model's accuracy, nor does it establish a
universal worst case over all dependence structures. It shows that marginal
performance and the number of agreeing judges do not determine the error of their
composition. A new model name or a different prompt is insufficient evidence of
independence.~\citep{exec24_scimethod_reference}

The bank-facing output should distinguish process completion, answer acceptance,
correct refusal, and unresolved cases. A confidence statement additionally needs
a defined event, population and task distribution, a sampling design, and a
justified construction evaluated on tasks not used for search or repair. An
independent causal oracle for generated models is one useful early tool. It can
check whether an agent preserves ambiguity or refuses an unidentified request,
instead of rewarding a fluent numerical answer. Evaluation on public observations
without intervention truth supports more limited checks of provenance,
reproducibility and assumptions. A model's own explanation can guide debugging,
but the perspective's concerns about self-explanation prevent treating it as a
faithful account of the model's internal mechanism.

\subsection{What public evidence can establish before bank experiments}

The product proposal should earn the expense of bank experiments through a staged
public and generated evidence portfolio. Generated structural models can supply
known intervention answers, nonidentification pairs, deliberate clock errors and
controlled measurement loss. Public empirical data can exercise extraction,
missingness, computation and reproducibility on realistic structures, with the
available causal truth stated separately. Open scientific tasks can test planning
and tool use, but their endpoints and safety constraints need not match banking.
This division lets inexpensive evidence test a concrete architectural proposition
without claiming that public data contain unavailable bank truth.

Transport remains a separate proposition. In a stipulated public environment,
contact and waiting yield net values USD 8 and USD 3 per account. A system can
select contact perfectly. In a second, explicitly different bank environment,
their values are USD $-4$ and USD 3. Applying the public choice incurs USD 7
regret relative to waiting. These are constructed utilities, not estimates. The
example isolates the need for a transport argument: good evaluation inside one
world does not establish which world the bank inhabits. Shared mechanisms,
covariates and measurement definitions may support a restricted bridge, but
those restrictions require evidence.~\citep{exec24_scimethod_reference}

For a first prototype, compare a fixed human-authored workflow, a retrieval-assisted
assistant, and a bounded planning assistant on the same frozen tasks and evidence.
The comparison should fix the allowed tools, retry budget, development/evaluation
split, answer and refusal rules, and costs before the final run. Include failures
to recover a paper, incorrectly interpreted clocks, ambiguous diagrams, unsupported
identification, malformed tool responses and stale results. A workflow that
correctly refuses an unanswerable request has a different achievement from one
that returns an accurate identified effect; both should be reported explicitly.
This is a proposed protocol, not an executed experiment or a claim of independent
acceptance. The number and mix of trials should follow the declared precision
and decision requirements, rather than an arbitrary success threshold.

\subsection{An auditable loop through the CLI and MCP}

The scientific loop motivates a concrete proposed interface. A question artifact
should bind the target, source identities and information clock. A diagram artifact
should carry alternative graphs and the provenance of constraints. An identification
request should return an estimand, an identified set, or a reasoned refusal. An
experiment proposal should bind the competing hypotheses and distinguishing
observations; an evaluation result should bind the exact task, methods and outputs.
CLI commands and MCP tools can expose these same typed artifacts, so that a
planner's change of wording cannot silently change the scientific target.

Acquisition, extraction, proof or numerical checking, and any execution have
different permissions and evidence duties. The proposal should keep hypothesis
content separate from permission to act. Before a tool invocation the host needs
to validate the requested operation, artifact identity and allowed context;
afterwards it needs to preserve the actual result and failure state. Retries
must retain earlier errors and cost, and resumed sessions must bind the same
frozen request. A local exact reference, an external package worker, and a
human-reviewed result are different execution modes. These contracts describe
required behavior; this source review has not implemented or qualified them.

Directly adopting a perspective as executable infrastructure is inappropriate:
it supplies an organizing argument, not a tested interface or algorithm. A useful
adoption is its separation of observation, hypothesis and experiment, accompanied
by the stricter causal and evidential contracts developed here. Porting an agent
framework is not justified by its appearance in the survey. The relevant package,
version, licensing, dependencies, termination and data boundaries need their own
primary review. A small explicit state machine remains a substantive alternative
to a general planner when the workflow is known; a planner may offer flexibility
when the next useful observation is genuinely unresolved. That flexibility must
be evaluated against its additional error, cost and review burden.

Finally, assume a fully counted exploratory workflow costs USD 240 for 40 attempts,
of which 30 complete their tools and eight deliver accepted scientific answers.
Its costs are USD 6 per attempt, USD 8 per tool completion and USD 30 per accepted
answer. Six correct refusals are reported separately; they are not silently added
to the answer denominator. All amounts and counts are stipulated. Recording
failed hypotheses and unsuccessful tool calls, consistent with the perspective's
case for learning from negative results, is necessary to estimate eventual
economics honestly. The contribution here is a testable proposal with explicit
information, translation, dependence and transport limits; native performance,
general algorithmic confidence and bank value remain unestablished.
\citep{exec24_scimethod_original,exec24_scimethod_reference}
